%=====================================================================
%  04_wheel.tex --- The desiccant wheel
%  Code of record: geothermal_cooling/dbhe_master.ipynb (v10.6), cell 12
%  (notebook cell index 26) for the model, cell 2 (index 6) for eta_F1/eta_F2.
%=====================================================================
\section{The desiccant wheel}
\label{sec:dw:main}

The wheel is the component that makes this plant a \emph{dehumidification}
plant rather than a heating plant, and it is where the geothermal heat is
finally spent. A rotary desiccant wheel is a porous matrix coated with a
sorbent --- silica gel here --- rotating slowly between two counter-flowing
air streams. In the \emph{process} sector it adsorbs moisture from the air
destined for the building; in the \emph{regeneration} sector, bathed in the
hot air leaving the coil, it desorbs that moisture and is readied for the
next revolution. Everything the well does upstream exists to keep that
regeneration air hot enough.

\subsection{Notation}
\label{sec:dw:notation}

Four air states matter. The third column gives the name carried in the
notebook, so that report and code can be read side by side.

\begin{center}
\begin{tabular}{@{}llll@{}}
\toprule
Station & Symbol & Notebook name & Origin \\
\midrule
Process inlet      & $(T_{\mathrm{pi}},w_{\mathrm{pi}})$ & \texttt{T\_p, w\_p}     & room return + fresh makeup \\
Process outlet     & $(T_{\mathrm{pd}},w_{\mathrm{pd}})$ & \texttt{T\_po, w\_po}   & dried air, feeds the cooler \\
Regeneration inlet & $(T_{\mathrm{ro}},w_{\mathrm{ri}})$ & \texttt{T\_reg, w\_reg} & coil outlet ($w$ unchanged: the coil is sensible) \\
Regeneration outlet& $(T_{\mathrm{rd}},w_{\mathrm{rd}})$ & \texttt{T\_rd, w\_rd}   & wheel exhaust, vented \\
\bottomrule
\end{tabular}
\end{center}

Dry-air mass flows are \mdp\ (process) and \mdr\ (regeneration), and in the
code $\mdr = 0.70\,\mdp$ through the configuration parameter
\texttt{reg\_pro\_ratio}. Humidity ratios $w$ are in
\si{\kilogram\per\kilogram} of dry air, and moist-air enthalpy is taken from
the same psychrometric block used everywhere else in the model,
\begin{equation}
h(T,w) = 1.006\,T + w\,(2501.0 + 1.86\,T)
\qquad [\si{\kilo\joule\per\kilogram}],\quad T\ \text{in \degC},
\label{eq:dw:enthalpy}
\end{equation}
with the inverse $T(h,w) = (h - 2501.0\,w)/(1.006 + 1.86\,w)$.

\subsection{Why a rotating matrix is hard}
\label{sec:dw:hard}

Inside the matrix two conservation laws hold at once and are strongly
coupled. Writing $z$ for distance along the flow, $\tau$ for time in the
rotation, $w$ and $h$ for the air humidity ratio and enthalpy, and $W$ and
$H$ for the \emph{sorbent} water loading and enthalpy,
\begin{align}
\text{moisture:}&\qquad
\frac{\partial w}{\partial z}+\frac{\partial W}{\partial \tau}=0,
\label{eq:dw:moisture_pde}\\
\text{energy:}&\qquad
\frac{\partial h}{\partial z}+\frac{\partial H}{\partial \tau}=0 .
\label{eq:dw:energy_pde}
\end{align}
The coupling runs through the sorption equilibrium. The loading $W=W(T,w)$
follows the isotherm, so it depends on \emph{both} air properties, and the
sorbent enthalpy $H=H(T,W)$ carries the heat of adsorption released when
water binds to the gel. The physical loop is short and vicious: adsorbing
water releases heat, the heat raises the local temperature, the higher
temperature shifts the equilibrium, and the shifted equilibrium changes how
much more water can be adsorbed. Temperature and humidity cannot be solved
one at a time.

Equations~\eqref{eq:dw:moisture_pde}--\eqref{eq:dw:energy_pde} form a coupled
non-linear hyperbolic system. Solving them directly means a spatially and
temporally resolved simulation of the matrix, repeated inside every step of a
twenty-year march that already carries a fixed-point loop around the well ---
unaffordable, and in any case dependent on matrix data (channel geometry,
sorbent mass, isotherm) that we do not have for a specific commercial wheel.

\subsection{The Banks/Close non-linear analogy}
\label{sec:dw:analogy}

The route out is the non-linear analogy method of Close and Banks, developed
for combined heat-and-mass regenerators and reduced to a usable form for
silica gel by Jurinak (1982). The idea is a change of variables. One seeks
two \emph{characteristic potentials} $F_1(T,w)$ and $F_2(T,w)$ such that the
coupled system becomes two \emph{independent} wave equations, each potential
being conserved along its own characteristic and travelling at its own speed
$\gamma_i$:
\begin{equation}
\frac{\partial F_i}{\partial z}
+\frac{1}{\gamma_i}\frac{\partial F_i}{\partial \tau}=0,
\qquad i=1,2 .
\label{eq:dw:waves}
\end{equation}

What this buys is decoupling. In $(F_1,F_2)$ coordinates the entangled
heat-and-mass problem becomes two separate single-variable transfer problems,
each characterised by one effectiveness exactly as a sensible heat exchanger
is characterised by one $\varepsilon$. The cost is that the potentials are not
universal: they are tied to the sorbent's isotherm, so the fits below are
silica-gel fits and nothing else.

\subsection{The two potentials, exactly as coded}
\label{sec:dw:potentials}

The notebook codes Jurinak's silica-gel potentials as
\begin{equation}
\boxed{\;
F_1(T,w) = -\,\frac{2865.0}{T^{1.490}} + 4.344\,w^{0.8624}
\;}
\label{eq:dw:F1}
\end{equation}
\begin{equation}
\boxed{\;
F_2(T,w) = \frac{T^{1.490}}{6360.0} - 1.127\,w^{0.07969}
\;}
\label{eq:dw:F2}
\end{equation}
which are, character for character, the two lines
\begin{center}
\texttt{def F1(T\_K, w): return -2865.0*T\_K**(-1.490) + 4.344*w**0.8624}\\
\texttt{def F2(T\_K, w): return\ \ T\_K**(1.490)/6360.0 - 1.127*w**0.07969}
\end{center}
of cell 12 of the model of record.

\paragraph{Units warning --- read this before using the formulas.}
Equations~\eqref{eq:dw:F1}--\eqref{eq:dw:F2} are \emph{dimensional} fits, not
dimensionless groups. $T$ \textbf{must} be in kelvin and $w$ \textbf{must} be
in \si{\kilogram\per\kilogram} of dry air, never \si{\gram\per\kilogram}.
Feeding \degC\ or \si{\gram\per\kilogram} produces numbers, no warning, and a
silently wrong wheel; with $w$ in \si{\gram\per\kilogram} the first term of
$F_2$ is unchanged while the second grows by a factor $1000^{0.07969}\approx
1.73$, which is enough to invert the direction of the answer. The notebook
respects this by converting at the single entry point of the wheel routine,
\texttt{Tp\_K, Tr\_K = T\_p+273.15, T\_reg+273.15}, and by carrying $w$ in
\si{\kilogram\per\kilogram} everywhere in the psychrometric block. Every
temperature that crosses the wheel boundary in either direction is in \degC;
kelvin exists only inside.

\subsection{What the potentials mean physically}
\label{sec:dw:meaning}

The two potentials have an interpretation that is also a way to test them.

Lines of constant $F_1$ are very nearly lines of constant \emph{enthalpy}.
$F_1$ is the fast potential, associated with the combined heat-and-mass wave
that sweeps quickly through the matrix; it is very nearly the
adiabatic-saturation invariant of the air, which is why holding $F_1$ fixed is
almost the same as holding $h$ fixed.

Lines of constant $F_2$ are very nearly lines of constant \emph{relative
humidity}. $F_2$ is the slow potential, associated with the sorption wave, and
it is the moisture-capacity coordinate: since $\partial F_2/\partial w<0$,
drying the air \emph{raises} $F_2$, and dragging the process state toward the
hot, low-\RH\ regeneration state is precisely dragging it up in $F_2$.

\paragraph{The self-check, and how the sign in $F_2$ was settled.}
Published transcriptions of Eq.~\eqref{eq:dw:F2} disagree on the sign of the
second term, so the notebook resolves it numerically and re-runs the check on
every execution. It evaluates $F_1$ along a constant-enthalpy line
($h = \SI{85}{\kilo\joule\per\kilogram}$) and $F_2$ along a constant relative
humidity line ($\RH = 0.40$):

\begin{center}
\begin{tabular}{@{}cccc@{}}
\toprule
$T$ [\degC] & $F_1$ at $h=\SI{85}{\kilo\joule\per\kilogram}$
            & $F_2$ at $\RH=0.40$ (coded, minus)
            & $F_2$ at $\RH=0.40$ (plus) \\
\midrule
25 & $-0.41814$ & $-0.00135$ & $+1.53061$ \\
45 & $-0.41621$ & $+0.00376$ & $+1.68083$ \\
55 & $-0.41908$ & $+0.00770$ & $+1.75639$ \\
\midrule
spread & $0.00287$ & $0.00905$ & $0.22578$ \\
\bottomrule
\end{tabular}
\label{tab:dw:selfcheck}
\end{center}

$F_1$ moves by $0.69\%$ of its own magnitude across a \SI{30}{\kelvin} span,
so the iso-enthalpy reading holds well. For $F_2$ the test is comparative
rather than absolute: the coded minus sign gives an iso-\RH\ line
\emph{twenty-five times flatter} than the plus sign would
($0.00905$ against $0.22578$), which settles the transcription without appeal
to any of the conflicting sources. This check is not decoration --- it is the
only defence the model has against a sign error that would otherwise be
invisible in the plant results.

\subsection{The effectiveness method}
\label{sec:dw:effectiveness}

Because the transformed problem~\eqref{eq:dw:waves} is decoupled, each
potential behaves like a single-stream transfer problem and can be assigned an
effectiveness. Define, exactly as $\varepsilon$ is defined for a sensible
exchanger --- actual change over the maximum change available ---
\begin{equation}
\eta_{Fi} \equiv
\frac{F_i^{\mathrm{pd}}-F_i^{\mathrm{pi}}}
     {F_i^{\mathrm{ri}}-F_i^{\mathrm{pi}}},
\qquad i=1,2,
\label{eq:dw:etadef}
\end{equation}
where $F_i^{\mathrm{ri}} \equiv F_i(T_{\mathrm{ro}},w_{\mathrm{ri}})$ is the
potential of the \emph{regeneration inlet}: the maximum change available to
the process stream is the full excursion from its own inlet to the state on
the other side of the matrix. Rearranged, this is the working form
coded in the notebook: the process state is dragged from its own inlet
towards the regeneration inlet, in $(F_1,F_2)$ space, by the prescribed
fractions,
\begin{equation}
\boxed{\;
F_i^{\mathrm{pd}} = F_i(T_{\mathrm{pi}},w_{\mathrm{pi}})
+\eta_{Fi}\Big[F_i(T_{\mathrm{ro}},w_{\mathrm{ri}})
-F_i(T_{\mathrm{pi}},w_{\mathrm{pi}})\Big],
\qquad i=1,2 .
\;}
\label{eq:dw:drag}
\end{equation}
The geometric picture is worth holding on to: the outlet sits a prescribed
fraction of the way along the segment joining the two inlets --- but that
segment is straight in $(F_1,F_2)$ coordinates, not on the psychrometric
chart, and it is the mapping back that makes the process path curve.

\paragraph{The values, and why they are asymmetric.}
The configuration cell prescribes
\begin{equation}
\eta_{F1}=0.10,\qquad \eta_{F2}=0.70 .
\label{eq:dw:etavals}
\end{equation}
The asymmetry is the design intent of a dehumidification wheel. A large
$\eta_{F2}$ means the wheel transfers most of the available sorption (that
is, relative-humidity) potential --- that is its job. A small $\eta_{F1}$
means it transfers little of the enthalpy potential, because transferring
enthalpy from the regeneration side would simply carry the regeneration heat
into the process stream, which the downstream cooler would then have to
remove again. A good wheel has $\eta_{F2}$ high and $\eta_{F1}$ low.

\paragraph{Standing flag.} These two numbers are literature-typical
placeholders for silica gel, and the configuration cell says so in a comment
(\texttt{<- FIT to wheel data}). They are constants only at fixed rotation
speed and balanced flow. Every quantitative statement in this report that
depends on the wheel inherits their uncertainty, and they must be fitted to
the catalogue or test data of a chosen wheel before any of it is claimed as a
prediction.

\subsection{Recovering the state: the 2-D Newton inversion}
\label{sec:dw:newton}

Equation~\eqref{eq:dw:drag} delivers the outlet \emph{potentials}, but the
rest of the plant needs the outlet \emph{state}
$(T_{\mathrm{pd}},w_{\mathrm{pd}})$. Since
Eqs.~\eqref{eq:dw:F1}--\eqref{eq:dw:F2} are non-linear in both arguments and
not analytically invertible, the code solves the $2\times2$ system
\begin{equation}
\mathbf{r}(T,w)=
\begin{pmatrix}
F_1(T,w)-F_1^{\mathrm{pd}}\\[2pt]
F_2(T,w)-F_2^{\mathrm{pd}}
\end{pmatrix}=\mathbf{0}
\label{eq:dw:residual}
\end{equation}
by Newton's method, with the Jacobian obtained by differentiating the fits
directly:
\begin{equation}
\mathbf{J}=
\begin{pmatrix}
\dfrac{\partial F_1}{\partial T} & \dfrac{\partial F_1}{\partial w}\\[10pt]
\dfrac{\partial F_2}{\partial T} & \dfrac{\partial F_2}{\partial w}
\end{pmatrix}
=
\begin{pmatrix}
+2865.0\times 1.490\;T^{-2.490} & 4.344\times 0.8624\;w^{-0.1376}\\[6pt]
\dfrac{1.490\,T^{0.490}}{6360.0} & -\,1.127\times 0.07969\;w^{-0.92031}
\end{pmatrix}.
\label{eq:dw:jacobian}
\end{equation}
Writing $\mathbf{J}=\begin{psmallmatrix}a&b\\c&d\end{psmallmatrix}$ with
$\det = ad-bc$, the update $\Delta\mathbf{x}=-\mathbf{J}^{-1}\mathbf{r}$ uses
the closed-form $2\times2$ inverse
$\mathbf{J}^{-1}=\tfrac{1}{\det}\begin{psmallmatrix}d&-b\\-c&a\end{psmallmatrix}$,
so that
\begin{equation}
\Delta T = -\,\frac{d\,r_1 - b\,r_2}{\det},
\qquad
\Delta w = -\,\frac{-c\,r_1 + a\,r_2}{\det},
\label{eq:dw:update}
\end{equation}
which is line for line what \texttt{\_invert\_F1F2} evaluates. The Jacobian
signs are the ones already discussed: $\partial F_1/\partial T>0$ and
$\partial F_2/\partial w<0$.

The iteration starts from the process inlet itself,
$(T,w)=(T_{\mathrm{pi}},w_{\mathrm{pi}})$ in kelvin, which is always a good
guess because the outlet is only a partial drag away from it. After each
update $w$ is floored, $w \leftarrow \max(w,10^{-7})$, to protect the
fractional powers $w^{-0.1376}$ and $w^{-0.92031}$ from a zero or negative
argument. The routine converts back on exit and returns $T$ in \degC.
Convergence is quadratic; the shipped self-check converges in five iterations.

\paragraph{What happens if it does not converge --- and this matters.}
The loop runs a fixed \texttt{for \_ in range(60)}. Its exit test,
$|r_1|<10^{-9}$ and $|r_2|<10^{-9}$, is applied to the residual computed
\emph{before} the current update, so on success the routine performs one
harmless extra step. On failure there is no branch at all: after sixty
iterations the routine simply \texttt{return}s the last iterate, with no
convergence flag, no exception, and no warning. A non-converged inversion is
therefore indistinguishable, downstream, from a converged one. The only mode
that raises rather than lies is a singular Jacobian, $\det=0$, which produces
a \texttt{ZeroDivisionError} on native floats. Adding a convergence flag that
propagates into the facility diagnostics is a small change and should be made
before the wheel is exercised far outside the conditions tested here.

\subsection{The regeneration outlet: conservation, not correlation}
\label{sec:dw:regen}

The effectiveness method gives only the process outlet. The regeneration
outlet comes from overall balances on the wheel --- which is exactly why the
process side must be closed before the regeneration side can be, and that
ordering propagates outward into the structure of the whole air-loop
solution.

Treat the wheel as \textbf{adiabatic} and at \textbf{cyclic steady state}: the
sorbent stores nothing net over a full revolution, and nothing is lost to the
surroundings. Then whatever moisture and enthalpy leave one stream must enter
the other. With
$\Delta w_{\mathrm{p}} \equiv w_{\mathrm{pi}} - w_{\mathrm{pd}}$ the moisture
removed from the process air (the quantity \texttt{dW\_process} returned by
the routine),
\begin{equation}
\mdp\,\Delta w_{\mathrm{p}} = \mdr\,(w_{\mathrm{rd}}-w_{\mathrm{ri}})
\quad\Longrightarrow\quad
\boxed{\;w_{\mathrm{rd}} = w_{\mathrm{ri}}
+ \frac{\mdp}{\mdr}\,\Delta w_{\mathrm{p}}\;}
\label{eq:dw:wrd}
\end{equation}
\begin{equation}
\mdp\,\big(h_{\mathrm{pi}}-h_{\mathrm{pd}}\big)
= \mdr\,\big(h_{\mathrm{rd}}-h_{\mathrm{ro}}\big)
\quad\Longrightarrow\quad
\boxed{\;h_{\mathrm{rd}} = h_{\mathrm{ro}}
+ \frac{\mdp}{\mdr}\,\big(h_{\mathrm{pi}}-h_{\mathrm{pd}}\big)\;}
\label{eq:dw:hrd}
\end{equation}
and finally $T_{\mathrm{rd}} = T(h_{\mathrm{rd}},w_{\mathrm{rd}})$ from the
inverse of Eq.~\eqref{eq:dw:enthalpy}.

\paragraph{Reading the signs.} When the wheel dries,
$\Delta w_{\mathrm{p}}>0$, so Eq.~\eqref{eq:dw:wrd} gives
$w_{\mathrm{rd}}>w_{\mathrm{ri}}$: the regeneration air picks the moisture up,
which is what desorption means. At the same time $\eta_{F1}$ is small but
positive, so the process air gains a little enthalpy,
$h_{\mathrm{pd}}>h_{\mathrm{pi}}$, and Eq.~\eqref{eq:dw:hrd} then gives
$h_{\mathrm{rd}}<h_{\mathrm{ro}}$: the regeneration air \emph{cools} as it
desorbs. Both directions are correct, and neither was imposed --- they fall
out of the balances once the process side is known.

\paragraph{Where conservation alone is not enough.} Equations
\eqref{eq:dw:wrd}--\eqref{eq:dw:hrd} carry no bound. They will return a
regeneration outlet at any temperature the arithmetic demands, including
states below the process inlet wet-bulb, if the flow ratio is pushed. At
$\mdr/\mdp=0.30$ with the self-check inlets, the balance returns
$T_{\mathrm{rd}}=\SI{-3.2}{\degreeCelsius}$ --- arithmetically consistent and
physically impossible. Nothing in the wheel routine flags this, because
nothing in the wheel routine knows about transfer limits. It is a reason to
keep $\mdr/\mdp$ near its configured value of $0.70$ unless the outlet state
is inspected.

\subsection{Why the flow ratio is a lever on one side only}
\label{sec:dw:lever}

The ratio $\mdp/\mdr$ appears in Eqs.~\eqref{eq:dw:wrd}--\eqref{eq:dw:hrd}
and nowhere else. On the \emph{regeneration} side it is a real lever: a small
regeneration flow concentrates the transferred moisture and enthalpy into a
small stream, giving a large rise in $w_{\mathrm{rd}}$ and a large drop in
$T_{\mathrm{rd}}$, while a large flow dilutes both. Combined with the coil's
temperature-versus-throughput trade-off, $\mdr/\mdp$ is one of the more
consequential design choices in the plant.

On the \emph{process} side it is not a lever at all. Equations
\eqref{eq:dw:drag}--\eqref{eq:dw:update} contain no mass flow of any kind.
The process outlet $(T_{\mathrm{pd}},w_{\mathrm{pd}})$ depends only on the two
inlet states and on $\eta_{F1},\eta_{F2}$. Doubling both air flows, or
halving the regeneration flow alone, leaves the process outlet bit-for-bit
identical; only the regeneration outlet moves. Numerically, with the
self-check inlets, $\mdr/\mdp \in \{0.30,\,0.70,\,2.00\}$ all return
$T_{\mathrm{pd}}=\SI{57.86}{\degreeCelsius}$ and
$w_{\mathrm{pd}}=\SI{12.33}{\gram\per\kilogram}$, while $T_{\mathrm{rd}}$
moves from $-3.2$ to $43.4$ to \SI{67.0}{\degreeCelsius}. Any conclusion
drawn from the model about process-side flow effects is an artefact of
something else in the plant, not of the wheel.

\subsection{The central consequence: latent becomes sensible}
\label{sec:dw:latent}

Because $\eta_{F1}\ll 1$, the process path is nearly iso-enthalpic --- $F_1$
barely moves, and constant $F_1$ is very nearly constant $h$. Along a
constant-enthalpy line, Eq.~\eqref{eq:dw:enthalpy} forces a reduction in $w$
to be paid for in $T$:
\begin{equation}
h=\text{const}=1.006\,T+w\,(2501+1.86\,T)
\quad\Longrightarrow\quad
\frac{\mathrm{d}T}{\mathrm{d}w}
\approx -\,\frac{2501}{1.006+1.86\,w} < 0 .
\label{eq:dw:latent2sensible}
\end{equation}
Numerically, removing \SI{1}{\gram\per\kilogram} at constant enthalpy heats
the air by about $2501/1.04 \approx \SI{2.4}{\kelvin}$.

This is not a defect of the wheel; it is thermodynamics. The latent heat
released when vapour is adsorbed does not vanish --- it reappears as sensible
heat in the very same air stream, plus whatever extra enthalpy the small
$\eta_{F1}$ drag imports from the hot regeneration side. \textbf{A desiccant
wheel is a latent-to-sensible converter.} It follows that the air-to-air
cooler downstream of the wheel is not an optional refinement or a
polish stage: without it the wheel delivers air that is drier than the room
but far hotter, and the building cannot be held at comfort. The wheel creates
the sensible load that the cooler --- and, in the modes that use it, the
chiller --- exists to reject.

\subsection{The regime diagnostic: when the wheel runs backwards}
\label{sec:dw:regime}

A desiccant wheel dries the process stream only while the regeneration air is
drier \emph{in relative-humidity terms} than the process air. Below that
crossover the drag in Eq.~\eqref{eq:dw:drag} points the wrong way, the wheel
humidifies the process air, and $\Delta w_{\mathrm{p}}$ simply goes negative.
Nothing breaks: the fixed-point loop still converges, the humidity controller
still reports its target met, and the reversal is invisible in the plant
results.

The notebook therefore instruments it. A module-level dictionary
\texttt{DW\_REGIME} counts every wheel evaluation and every reversed one,
records the worst (most negative) $\Delta w_{\mathrm{p}}$, and raises a single
\texttt{UserWarning} on the first reversal quoting both inlet states, both
relative humidities, and the \RH\ gap that must stay positive; at the end of a
run \texttt{dw\_regime\_report()} prints the reversed fraction and the worst
excursion. This is a reporting device, not a correction: the model still
returns the reversed state and the plant still runs on it. It exists so that a
reversal is never silent.

\subsection{Worked example: the shipped self-check}
\label{sec:dw:example}

The last lines of cell 12 exercise the wheel at a Gulf-Coast-like process
inlet, \SI{32}{\degreeCelsius} at $\RH=0.68$, against a regeneration inlet at
\SI{80}{\degreeCelsius} carrying the same absolute humidity, at the
configured flow ratio $\mdr/\mdp=0.70$.

\begin{center}
\begin{tabular}{@{}lccccl@{}}
\toprule
Stream & $T$ [\degC] & $w$ [\si{\gram\per\kilogram}]
       & $h$ [\si{\kilo\joule\per\kilogram}] & \RH\ [\%] & \\
\midrule
Process inlet  (pi) & 32.0 & 20.48 & \phantom{0}84.6 & 68.0 & room return + makeup \\
Process outlet (pd) & 57.9 & 12.33 & \phantom{0}90.4 & 10.9 & \textbf{dried \SI{8.15}{\gram\per\kilogram}, heated \SI{25.9}{\kelvin}} \\
Regen inlet    (ro) & 80.0 & 20.48 & 134.7 & \phantom{0}6.7 & coil outlet \\
Regen outlet   (rd) & 43.4 & 32.11 & 126.5 & --- & \textbf{cooler and wetter} \\
\bottomrule
\end{tabular}
\label{tab:dw:example}
\end{center}

Check the moisture balance, Eq.~\eqref{eq:dw:wrd}:
\begin{equation*}
w_{\mathrm{rd}} = 20.48 + \frac{1}{0.70}\,(20.48-12.33)
= 20.48 + 1.4286\times 8.15 = \SI{32.11}{\gram\per\kilogram}.\quad\checkmark
\end{equation*}

Two readings. The process enthalpy rise is only
$90.4-84.6=\SI{5.7}{\kilo\joule\per\kilogram}$ out of an available
$134.7-84.6=\SI{50.1}{\kilo\joule\per\kilogram}$ --- the $\eta_{F1}=0.10$ drag
doing exactly what it was set to do. And the temperature rise of
\SI{25.9}{\kelvin} exceeds the pure iso-enthalpic estimate $2.4\times
8.15\approx\SI{19.6}{\kelvin}$ by just that imported enthalpy. The
regeneration stream mirrors it, leaving \SI{36.6}{\kelvin} colder and
\SI{11.6}{\gram\per\kilogram} wetter than it entered.

\subsection{Limitations: the formulation carries no size}
\label{sec:dw:limitations}

This subsection is the one to read before quoting any wheel result.

\paragraph{There is no wheel in the model.} Equations
\eqref{eq:dw:drag}--\eqref{eq:dw:update} contain no diameter, no depth, no
desiccant mass, no channel geometry, no transfer area, no rotation speed and
no NTU. The entire hardware description of the wheel is the pair of numbers
$(\eta_{F1},\eta_{F2}) = (0.10,\,0.70)$. Everything the model knows about the
component is compressed into those two constants, and they were prescribed,
not derived. Three consequences follow, and none is subtle.

\emph{The process outlet is independent of both mass flows.} Shown
numerically in Section~\ref{sec:dw:lever}: the flows enter only the
regeneration balances. In a real wheel the process outlet depends strongly on
face velocity, because face velocity sets residence time and therefore the
number of transfer units. The model cannot see this.

\emph{The wheel cannot be resized.} There is no parameter to change that
makes the wheel bigger or smaller. Scaling the plant --- changing
\texttt{N\_eggs}, and with it \mdp\ --- rescales every duty in the facility
but leaves the wheel's performance per kilogram of air exactly where it was.
A wheel sized for the design flow and a wheel half that size are the same
object inside this model.

\emph{Part-load and rotation-speed behaviour are outside the model's
vocabulary.} Wheel rotation speed is the primary control of a real desiccant
wheel: too slow and the matrix saturates before it returns to regeneration,
too fast and it carries hot regeneration air into the process stream as
parasitic heat, and the optimum shifts with load. None of this exists here.
Turning down the plant changes the inlet states, and the wheel responds to
those, but it responds with the same two effectiveness values it uses at full
load. Any statement in this report about wheel sizing, rotation speed,
face-velocity selection or part-load degradation is outside what the model can
represent, and should be read as an assumption or as a placeholder, never as
a result.

\paragraph{The remedy, and why it is deferred.} The analogy method itself
supplies the missing link. Once the problem is decoupled into the two waves of
Eq.~\eqref{eq:dw:waves}, each wave admits a regenerator-style effectiveness
correlation of the form $\eta_{Fi} = \eta_{Fi}(\mathrm{NTU}_i,
C_{r,i}^{*})$, where the number of transfer units follows from the transfer
coefficient and the transfer area, $\mathrm{NTU}_i \sim \UA_i/(\md\,c_p)_i$,
and the matrix capacity ratio follows from the desiccant mass and the
rotation speed. Supplying wheel geometry, sorbent mass and rotation speed
would then \emph{produce} $\eta_{F1}$ and $\eta_{F2}$ instead of taking them
as input, and every limitation listed above would lift at once: face-velocity
dependence, resizing and part-load behaviour would all follow.

That extension is deferred. It needs geometric and sorbent data for a specific
commercial wheel, which we do not have, and fitting it to a catalogue is a
separate exercise from the plant-level architecture this report documents.
Until it is done, $\eta_{F1}$ and $\eta_{F2}$ remain what the configuration
cell calls them: numbers to be fitted to wheel data.
